% Consolidated external validation paper: T1--T12 (see EPIC_016/017/018 specs in specs/COMPLETE/).
% Normative: EPIC_016_WV2A; EPIC_017_EV1 (T4--T7); EPIC_018_XC1 (T8--T12).
\documentclass[11pt]{article}
\input{../suite_preamble}
\usepackage{longtable}
\usepackage{listings}
\lstdefinelanguage{lean}{
  keywords={def, theorem, lemma, structure, class, instance, where, by, exact, rfl, simp, intro, have, let, refine, rcases, obtain, cases, fun},
  keywordstyle=\bfseries,
  ndkeywords={Prop, Type, Sort, Nat},
  ndkeywordstyle=\bfseries,
  sensitive=true,
  comment=[l]{--},
  morecomment=[s]{/-}{-/},
  commentstyle=\itshape,
  stringstyle=\ttfamily,
  basicstyle=\ttfamily\small,
  breaklines=true,
  showstringspaces=false
}
\input{../suite_macros}

\newcommand{\IC}{\ensuremath{\mathsf{IC}}}
\newcommand{\Fib}{\mathrm{Fib}}

\title{External Validation of a Positive-Closure Proof Architecture\\[0.35em]
\large A Multi-Tranche Survey: Twelve Benchmark Dimensions in Lean}
\author{Nova Spivack}
\date{2026}

\begin{document}

\maketitle

\begin{abstract}
The Infinity Compression (\IC{}) formalization develops a layered positive-closure proof
architecture: collapse to a bare certificate, a typed forgetful map, fiber structure over
that certificate, and, where appropriate, canonical witnesses or sections. A natural
methodological concern is that this architecture may be endogenous to the flagship
NEMS-style case rather than a reusable proof pattern. This paper presents a
\textbf{consolidated machine-checked external validation}: twelve independent tranches of
ordinary mathematics, each formalized in Lean without importing the \IC{} spine, each
comparing a \emph{flat baseline} (the obvious Mathlib-shaped existence or representability
statement) with an \emph{architecture-guided} theorem bundle (nonempty fibers, surjectivity
of the forgetful map, \texttt{HasRightInverse}, an identified \texttt{RightInverse} section,
and a uniform canonical witness in each fiber), \emph{except} tranche \textbf{T12}, which
records an explicit \textbf{MODERATE} / control pattern for a \emph{non-surjective} collapse.

Tranche \textbf{T1} is the setoid quotient benchmark (collapse by an equivalence relation;
baseline \texttt{Quotient.exists\_rep}). Tranche \textbf{T2} is the Cartesian product with
first projection (many enriched pairs over one bare image---\emph{multi-route} without
quotienting). Tranche \textbf{T3} is the dependent sum $\Sigma_{b:B} E\,b$ with
\texttt{Sigma.fst} (bundle projection; dependent generalization of T2). Tranches
\textbf{T4}--\textbf{T7} cover sum discriminants, subtype projection, orbit quotients, and
ideal quotients. Tranches \textbf{T8}--\textbf{T11} add pushout glueing, set membership as
classifier, localization via fraction pairs, and pullbacks (sheaf-style fiber products).
Tranche \textbf{T12} is \texttt{Nat.succ}: fibers are nonempty exactly off zero, with no
global right inverse. The contribution is
not a catalog of new classical theorems; it is a demonstration that the same architectural
lens organizes \emph{twelve} external families and yields strengthened, reusable proof
interfaces where the rubric applies. Under the project's validation rubric, tranches T1--T11
support a \textbf{STRONG} verdict; T12 documents a complementary \textbf{MODERATE} boundary
case. We situate
the work in the NEMS / APS / IC / Strata stack and state explicitly that APS is not tested
by these tranches.
\end{abstract}

\tableofcontents
\newpage

\section{Introduction}\label{sec:intro}

The Infinity Compression (\IC{}) program develops a positive-closure proof architecture
organized around four layers: a collapsed bare layer, a projection from enriched to bare
structure, fibers over bare points, and, when available, canonical witnesses or sections.
In the flagship \IC{} formalization this architecture is load-bearing: it organizes the
distinction between theorem-level certification and richer reflective realization.
\textbf{The central claim of this paper is that the positive-closure architecture is a
reusable proof pattern---not merely a custom narrative for the flagship case---and that
claim is supported by multiple independent external benchmarks.}

A single quotient formalization already begins to answer the worry that the architecture
might be idiosyncratic; this document \emph{consolidates} that quotient tranche with eleven
additional standard constructions---products, dependent sums, coproduct discriminants,
subtypes, orbit and ideal quotients, pushouts, classifiers, localization pairs, pullbacks,
and a non-surjective control---so that the
evidence base is explicitly multi-family. The result is one self-contained validation paper:
readers need not cross-reference a separate ``survey'' manuscript for the product and bundle
stories.

\paragraph{Scope.}
We do not claim a universal theorem that every closure problem in mathematics must be
formalized in \IC{}'s types. We claim that, in Lean, on twelve concrete external families,
the architecture exports cleanly, yields structured theorem bundles beyond flat baselines,
and aligns with the flagship forgetful / fiber / section narrative~\cite{SpivackIC}.

\section{The positive-closure architecture as a reusable schema}\label{sec:schema}

Let $\pi$ denote a forgetful projection from an \emph{enriched} carrier to a \emph{bare}
certificate space (a collapsed layer where theorem-level records live). Write $E$ for an
enriched carrier and $B$ for a bare layer. The schema is:
\begin{equation}\label{eq:schema}
E \xrightarrow{\;\pi\;} B,
\qquad
\Fib(b) \coloneqq \pi^{-1}(b)\ \text{ for } b\in B,
\qquad
\text{optionally } s : B \to E \text{ with } \pi \circ s = \mathrm{id}_B.
\end{equation}
Here $\Fib(b)$ is the fiber over a bare point; a section $s$ is a global choice of preimage.

Each tranche below instantiates \eqref{eq:schema} with different $E$, $B$, and $\pi$. The
quotient tranche uses Mathlib's quotient map and \texttt{Quotient.out}; the product tranche
uses first projection and a parametric section $x \mapsto (x,y_0)$; the bundle tranche uses
\texttt{Sigma.fst} and a \texttt{noncomputable} section from classical choice when
$\forall b,\ \mathrm{Nonempty}(E\,b)$.

\section{Unified methodology: baseline vs.\ architecture-guided bundle}\label{sec:method}

In every tranche we record:
\begin{itemize}[leftmargin=2em]
  \item a \textbf{baseline} statement profile matching how Mathlib users typically see the
    mathematics (existence of representatives, trivial points in fibers, etc.);
  \item an \textbf{architecture-guided bundle}: named bare layer, named projection, fiber
    type as a subtype or structured copy, surjectivity of $\pi$, \texttt{HasRightInverse},
    a concrete \texttt{RightInverse} witness, and a \texttt{canonical\ldots Witness} in
    each fiber.
\end{itemize}
The methodological gain is \emph{interface structure}: the bundle is a reusable template
parallel to the flagship story, not a claim of new classical content.

\section{Tranche T1: Setoid quotients}\label{sec:t1}

\subsection{Why quotients}\label{sec:t1-why}

This tranche is a positive-closure benchmark in the sense relevant to \IC{}: it exhibits
collapse, fibers, and a uniform section-like recovery mechanism (\texttt{Quotient.out}), not
only an existential shadow.

The quotient construction is standard mathematics (see, e.g., \cite{MacLane}); it has a
genuine collapse map from carriers to quotient points; it admits a classical baseline
(\texttt{Quotient.exists\_rep}); and it supports a section story via \texttt{Quotient.out}.
The benchmark imports only \texttt{Mathlib.Data.Setoid.Basic} and
\texttt{Mathlib.Logic.Function.Basic}.

\subsection{Lean formalization and main bundle}\label{sec:t1-lean}

All definitions and theorems live in
\texttt{InfinityCompression/Validation/QuotientFiberBenchmark.lean} (namespace
\texttt{InfinityCompression.Validation}). Layering: \texttt{BareQuotient},
\texttt{EnrichedCarrier}, \texttt{forgetToQuotient}, \texttt{QuotientFiber},
\texttt{canonicalQuotientFiberWitness}.

\begin{table}[htbp]
\centering
\footnotesize
\begin{tabular}{@{}>{\raggedright\arraybackslash}p{0.34\textwidth}>{\raggedright\arraybackslash}p{0.36\textwidth}>{\raggedright\arraybackslash}p{0.24\textwidth}@{}}
\toprule
\textbf{Lean name} & \textbf{Meaning} & \textbf{Baseline vs.\ architecture} \\
\midrule
\path{quotient_fiber_nonempty} &
each quotient point has a witness in its \texttt{QuotientFiber} &
baseline-plus (typed fiber nonempty) \\
\path{forgetToQuotient_surjective} &
\texttt{forgetToQuotient} hits every bare point &
architecture \\
\path{forgetToQuotient_hasRightInverse} &
projection admits a global right inverse &
architecture \\
\path{quotientOut_rightInverse_forgetToQuotient} &
\texttt{Quotient.out} is a \texttt{RightInverse} of \texttt{forgetToQuotient} &
architecture \\
\path{canonicalQuotientFiberWitness} &
uniform canonical witness in each fiber &
architecture \\
\bottomrule
\end{tabular}
\caption{T1 (quotient): main named results in the architecture-guided bundle.}
\label{tab:t1}
\end{table}

\subsection{What the architectural packaging adds}\label{sec:t1-pack}

Classically, every quotient class has \emph{some} representative (\texttt{Quotient.exists\_rep}).
The architecture-guided view reorganizes the same fact under \eqref{eq:schema}: surjectivity,
\texttt{HasRightInverse}, identification of \texttt{Quotient.out} as \texttt{RightInverse}, and
\texttt{canonicalQuotientFiberWitness}. Methodologically, this upgrades a pointwise existential
into a named, reusable interface---the same move the positive-closure architecture is designed
to expose internally.

\section{Tranche T2: Cartesian product and first projection}\label{sec:t2}

\subsection{Mathematical story}

Let $X$ and $Y$ be types. The first projection $\pi : X \times Y \to X$ forgets the second
coordinate. For fixed nonempty $Y$, many pairs $(x,y_1)$ and $(x,y_2)$ share the same bare
image $x$: this is \emph{multi-route over one bare certificate} without any quotient---a
clean stress test for ``fiber + section'' language distinct from equivalence collapsing.

\subsection{Baseline vs.\ bundle}

\textbf{Baseline:} once $Y$ is inhabited, each $x : X$ has some pair mapping to it (e.g.\
$(x,y_0)$ for a chosen $y_0$). \textbf{Architecture bundle:} we package
\texttt{ProductFiberAt X Y x}, surjectivity of \texttt{forgetFirst}, \texttt{HasRightInverse},
the parametric section \texttt{productSection y0}, \texttt{RightInverse} compatibility, and
\texttt{canonicalProductFiberWitness}. The Lean module is
\texttt{ProductProjectionFiberBenchmark.lean}; it imports only
\texttt{Mathlib.Logic.Function.Basic}.

\begin{table}[htbp]
\centering
\footnotesize
\begin{tabular}{@{}>{\raggedright\arraybackslash}p{0.36\textwidth}>{\raggedright\arraybackslash}p{0.56\textwidth}@{}}
\toprule
\textbf{Lean name} & \textbf{One-line reading} \\
\midrule
\path{forgetFirst} & first projection $X \times Y \to X$. \\
\path{ProductFiberAt X Y x} & fiber of \path{forgetFirst} over \path{x} (\path{Y} explicit in the name). \\
\path{product_fiber_nonempty} & each bare \path{x} has a nonempty fiber (given \path{y0 : Y}). \\
\path{forgetFirst_surjective} & \path{forgetFirst} is surjective (given \path{y0}). \\
\path{forgetFirst_hasRightInverse} & surjectivity $\Rightarrow$ \texttt{HasRightInverse}. \\
\path{productSection_rightInverse_forgetFirst} & \path{productSection y0} is \texttt{RightInverse} to \path{forgetFirst}. \\
\path{canonicalProductFiberWitness} & distinguished point \path{⟨(x, y0), rfl⟩} in each fiber. \\
\bottomrule
\end{tabular}
\caption{T2 (product projection): architecture-guided bundle.}
\label{tab:t2}
\end{table}

\section{Tranche T3: Dependent sums and the bundle projection}\label{sec:t3}

\subsection{Mathematical story}

For $E : B \to \mathrm{Type}$, the dependent sum $\Sigma_{b:B} E\,b$ is the standard
\emph{bundle} total space; \texttt{Sigma.fst} projects to the base. The fiber over
$b : B$ is (equivalent to) $E\,b$. This generalizes T2: product is the special case of a
constant family $E\,b = Y$.

\subsection{Hypothesis and classical section}

We assume $\forall b : B,\ \mathrm{Nonempty}(E\,b)$ so that \texttt{Sigma.fst} is surjective.
The global section \path{sigmaSection} uses \path{Classical.choose} per fiber; the
module is \path{noncomputable} where required. Baseline: each base point has \emph{some}
fiber point. Bundle: \path{SigmaFiber}, surjectivity, \path{HasRightInverse},
\path{sigmaSection_rightInverse_sigma_fst}, \path{canonicalSigmaFiberWitness}.

\begin{table}[htbp]
\centering
\footnotesize
\begin{tabular}{@{}>{\raggedright\arraybackslash}p{0.36\textwidth}>{\raggedright\arraybackslash}p{0.56\textwidth}@{}}
\toprule
\textbf{Lean name} & \textbf{One-line reading} \\
\midrule
\path{SigmaFiber E b} & points in $\Sigma\, E$ lying over \path{b}. \\
\path{sigma_fiber_nonempty} & nonempty fibers from $\forall b,\ \mathrm{Nonempty}(E\,b)$. \\
\path{sigma_fst_surjective} & \texttt{Sigma.fst} surjective. \\
\path{sigma_fst_hasRightInverse} & \texttt{HasRightInverse} for \texttt{Sigma.fst}. \\
\path{chooseFiber} / \path{sigmaSection} & classical section of the bundle. \\
\path{sigmaSection_rightInverse_sigma_fst} & \texttt{RightInverse} law. \\
\path{canonicalSigmaFiberWitness} & uniform witness in each \texttt{SigmaFiber}. \\
\bottomrule
\end{tabular}
\caption{T3 (dependent sum): architecture-guided bundle (\texttt{SigmaFiberBenchmark.lean}).}
\label{tab:t3}
\end{table}

\section{Tranche T4: Sum types and discriminant collapse}\label{sec:t4}

The coproduct \texttt{Sum A B} carries a natural \emph{tag-forgetting} map to \texttt{Bool}
(\texttt{forgetSumIsLeft}), parallel in spirit to T2's projection but for coproducts rather
than products: many \texttt{inl} values share the same bare tag. The module is
\texttt{SumCoproductFiberBenchmark.lean}. Baseline: both sides inhabited to hit both booleans.
Bundle: \texttt{SumFiberAt}, surjectivity, \texttt{sumSection}, \texttt{RightInverse} law,
\texttt{canonicalSumFiberWitness}.

\begin{table}[htbp]
\centering
\footnotesize
\begin{tabular}{@{}>{\raggedright\arraybackslash}p{0.36\textwidth}>{\raggedright\arraybackslash}p{0.56\textwidth}@{}}
\toprule
\textbf{Lean name} & \textbf{One-line reading} \\
\midrule
\path{forgetSumIsLeft} & collapse \texttt{Sum A B} to a \texttt{Bool} side tag. \\
\path{SumFiberAt A B b} & fiber over a tag \texttt{b : Bool}. \\
\path{sumSection} & section \texttt{Bool → Sum A B} at fixed witnesses. \\
\path{forgetSumIsLeft_surjective} & surjectivity (given \texttt{a0}, \texttt{b0}). \\
\path{sumSection_rightInverse_forgetSumIsLeft} & \texttt{RightInverse} law. \\
\path{canonicalSumFiberWitness} & canonical witness in each \texttt{SumFiberAt}. \\
\bottomrule
\end{tabular}
\caption{T4 (sum / coproduct): architecture-guided bundle.}
\label{tab:t4}
\end{table}

\section{Tranche T5: Subtype projection}\label{sec:t5}

For a predicate \texttt{p} on a type $\alpha$, \texttt{Subtype.val} forgets proofs. Under the hypothesis
$\forall a,\ p\ a$, \texttt{Subtype.val} is surjective onto $\alpha$; fibers are
singletons. Module: \texttt{SubtypeValFiberBenchmark.lean}.

\begin{table}[htbp]
\centering
\footnotesize
\begin{tabular}{@{}>{\raggedright\arraybackslash}p{0.36\textwidth}>{\raggedright\arraybackslash}p{0.56\textwidth}@{}}
\toprule
\textbf{Lean name} & \textbf{One-line reading} \\
\midrule
\path{SubtypeFiberAt} & fiber of \texttt{Subtype.val} over \texttt{a :} $\alpha$. \\
\path{subtype_val_surjective} & surjectivity from $\forall a,\ p\ a$. \\
\path{subtypeSection} & section sending \texttt{a} to \texttt{<a, h a>}. \\
\path{subtypeSection_rightInverse_subtypeVal} & \texttt{RightInverse} law. \\
\path{canonicalSubtypeFiberWitness} & canonical witness in each fiber. \\
\bottomrule
\end{tabular}
\caption{T5 (subtype): architecture-guided bundle.}
\label{tab:t5}
\end{table}

\section{Tranche T6: Orbit quotient from a group action}\label{sec:t6}

Given \texttt{MulAction G X}, the orbit relation \texttt{same\_orbit} yields a
\texttt{Setoid} and quotient \texttt{BareOrbitQuotient}; the collapse map is
\texttt{to\_orbit\_quotient}. Module: \texttt{OrbitRelationQuotientFiberBenchmark.lean}.

\begin{table}[htbp]
\centering
\footnotesize
\begin{tabular}{@{}>{\raggedright\arraybackslash}p{0.36\textwidth}>{\raggedright\arraybackslash}p{0.56\textwidth}@{}}
\toprule
\textbf{Lean name} & \textbf{One-line reading} \\
\midrule
\texttt{same\_orbit} & orbit relation $\exists g,\ g \cdot x = y$. \\
\texttt{grp\_action\_orbit\_setoid} & \texttt{Setoid} for the orbit relation. \\
\texttt{to\_orbit\_quotient} & \texttt{Quotient.mk} from \texttt{X} to orbit space. \\
\texttt{to\_orbit\_quotient\_surjective} & surjectivity. \\
\texttt{quotient\_out\_rightInverse\_to\_orbit\_quotient} & \texttt{Quotient.out} as \texttt{RightInverse}. \\
\texttt{canonical\_orbit\_quotient\_fiber\_witness} & canonical witness in each fiber. \\
\bottomrule
\end{tabular}
\caption{T6 (orbit quotient): architecture-guided bundle.}
\label{tab:t6}
\end{table}

\section{Tranche T7: Ideal quotient}\label{sec:t7}

For a commutative ring \texttt{R} and ideal \texttt{I}, the map
\texttt{Ideal.Quotient.mk I} collapses to the quotient $R/I$. This is the standard algebraic
quotient parallel to T1's setoid quotient. Module: \texttt{IdealQuotientFiberBenchmark.lean}.

\begin{table}[htbp]
\centering
\footnotesize
\begin{tabular}{@{}>{\raggedright\arraybackslash}p{0.36\textwidth}>{\raggedright\arraybackslash}p{0.56\textwidth}@{}}
\toprule
\textbf{Lean name} & \textbf{One-line reading} \\
\midrule
\path{forgetToIdealQuotient} & \texttt{Ideal.Quotient.mk I}. \\
\path{IdealQuotientFiber} & fiber over a bare coset in $R/I$. \\
\path{forgetToIdealQuotient_surjective} & surjectivity (via \texttt{Quotient.exists\_rep}). \\
\path{idealQuotientOut_rightInverse_forgetToIdealQuotient} & \texttt{Quotient.out} as \texttt{RightInverse}. \\
\path{canonicalIdealQuotientFiberWitness} & canonical witness in each fiber. \\
\bottomrule
\end{tabular}
\caption{T7 (ideal quotient): architecture-guided bundle.}
\label{tab:t7}
\end{table}

\section{Tranche T8: Pushout / colimit glueing}\label{sec:t8}

Glueing \texttt{inl a0} with \texttt{inr b0} in \texttt{Sum A B} is presented as the equivalence
generated by \texttt{Relation.EqvGen} on explicit symmetric edges (\texttt{PushoutCoequalizerFiberBenchmark.lean}).

\begin{table}[htbp]
\centering
\footnotesize
\begin{tabular}{@{}>{\raggedright\arraybackslash}p{0.36\textwidth}>{\raggedright\arraybackslash}p{0.56\textwidth}@{}}
\toprule
\textbf{Lean name} & \textbf{One-line reading} \\
\midrule
\texttt{pushoutRel} & generating edges for glueing. \\
\texttt{forgetToPushout} & \texttt{Quotient.mk} to the pushout type. \\
\texttt{forgetToPushout\_surjective} & surjectivity. \\
\texttt{quotientOut\_rightInverse\_forgetToPushout} & \texttt{Quotient.out} as \texttt{RightInverse}. \\
\texttt{canonicalPushoutFiberWitness} & canonical witness in each fiber. \\
\bottomrule
\end{tabular}
\caption{T8 (pushout): architecture-guided bundle.}
\label{tab:t8}
\end{table}

\section{Tranche T9: Membership classifier}\label{sec:t9}

Subsets as maps to \texttt{Prop} and \texttt{Subtype} of membership in \texttt{S}: \texttt{SetClassifierFiberBenchmark.lean}
packages membership fibers and global sections when every \texttt{x} lies in \texttt{S}.

\begin{table}[htbp]
\centering
\footnotesize
\begin{tabular}{@{}>{\raggedright\arraybackslash}p{0.36\textwidth}>{\raggedright\arraybackslash}p{0.56\textwidth}@{}}
\toprule
\textbf{Lean name} & \textbf{One-line reading} \\
\midrule
\texttt{membershipClassifier} & characteristic predicate \texttt{x in S}. \\
\texttt{ClassifierFiber} & fiber over \texttt{x} for \texttt{Subtype.val}. \\
\texttt{classifier\_fiber\_nonempty\_of\_mem} & nonempty fiber when \texttt{x} lies in \texttt{S}. \\
\texttt{subtypeVal\_surjectiveOn} & surjectivity of \texttt{Subtype.val} under global membership. \\
\texttt{classifierSection\_rightInverse\_subtypeVal} & \texttt{RightInverse} for the section. \\
\bottomrule
\end{tabular}
\caption{T9 (classifier / subset): architecture-guided bundle.}
\label{tab:t9}
\end{table}

\section{Tranche T10: Localization pairs}\label{sec:t10}

Fraction presentation \texttt{mk'} from \texttt{R × M} is surjective; \texttt{LocalizationPairFiberBenchmark.lean}
uses \texttt{IsLocalization.sec} as canonical preimage data.

\begin{table}[htbp]
\centering
\footnotesize
\begin{tabular}{@{}>{\raggedright\arraybackslash}p{0.36\textwidth}>{\raggedright\arraybackslash}p{0.56\textwidth}@{}}
\toprule
\textbf{Lean name} & \textbf{One-line reading} \\
\midrule
\texttt{forgetPairToLocalization} & \texttt{IsLocalization.mk'} on pairs. \\
\texttt{forgetPairToLocalization\_surjective} & pair-map surjectivity. \\
\texttt{sec\_rightInverse\_forgetPairToLocalization} & \texttt{IsLocalization.sec} as \texttt{RightInverse}. \\
\texttt{canonicalLocalizationPairFiberWitness} & canonical pair in each fiber. \\
\bottomrule
\end{tabular}
\caption{T10 (localization pairs): architecture-guided bundle.}
\label{tab:t10}
\end{table}

\section{Tranche T11: Pullback / fiber product}\label{sec:t11}

Given \texttt{f : A → C}, \texttt{g : B → C}, the subtype of \texttt{A × B} with \texttt{f a = g b}
is the pullback; \texttt{pullbackFst} projects to \texttt{A} (\texttt{PullbackFiberBenchmark.lean}).

\begin{table}[htbp]
\centering
\footnotesize
\begin{tabular}{@{}>{\raggedright\arraybackslash}p{0.36\textwidth}>{\raggedright\arraybackslash}p{0.56\textwidth}@{}}
\toprule
\textbf{Lean name} & \textbf{One-line reading} \\
\midrule
\texttt{PullbackType} & fiber-product subtype. \\
\texttt{pullbackFst} & first projection (renamed to avoid product-tranche clash). \\
\texttt{pullbackFst\_surjective} & surjectivity assuming overlap (for all \texttt{a} some \texttt{b} matches). \\
\texttt{pullbackSection\_rightInverse\_pullbackFst} & chosen section is \texttt{RightInverse}. \\
\bottomrule
\end{tabular}
\caption{T11 (pullback): architecture-guided bundle.}
\label{tab:t11}
\end{table}

\section{Tranche T12: Non-surjective collapse (MODERATE)}\label{sec:t12}

\texttt{Nat.succ} is not surjective; fibers over \texttt{n > 0} are nonempty singletons.
\texttt{NonSurjectiveSuccFiberBenchmark.lean} records \texttt{¬ HasRightInverse succ} and
\texttt{pred\_leftInverse\_succ}.

\begin{table}[htbp]
\centering
\footnotesize
\begin{tabular}{@{}>{\raggedright\arraybackslash}p{0.36\textwidth}>{\raggedright\arraybackslash}p{0.56\textwidth}@{}}
\toprule
\textbf{Lean name} & \textbf{One-line reading} \\
\midrule
\texttt{succFiber} & fiber of \texttt{succ} over \texttt{n}. \\
\texttt{succFiber\_nonempty\_iff} & nonempty iff \texttt{n} is not zero. \\
\texttt{nat\_succ\_not\_surjective} & obstruction to global section. \\
\texttt{pred\_leftInverse\_succ} & standard left-inverse identity. \\
\bottomrule
\end{tabular}
\caption{T12 (MODERATE control): \texttt{Nat.succ}.}
\label{tab:t12}
\end{table}

\section{Comparative synthesis across tranches}\label{sec:synth}

\begin{table}[htbp]
\centering
\footnotesize
\begin{tabular}{@{}l>{\raggedright\arraybackslash}p{0.22\textwidth}>{\raggedright\arraybackslash}p{0.22\textwidth}>{\raggedright\arraybackslash}p{0.26\textwidth}>{\raggedright\arraybackslash}p{0.18\textwidth}@{}}
\toprule
\textbf{Tr.} & \textbf{Bare $B$} & \textbf{Projection $\pi$} & \textbf{What varies in the fiber} & \textbf{Section flavor} \\
\midrule
T1 & quotient & \texttt{forgetToQuotient} / \texttt{Quotient.mk} & equivalence classes & \texttt{Quotient.out} (canonical) \\
T2 & $X$ & \texttt{forgetFirst} & second coordinate $y : Y$ & \texttt{productSection y0} (parametric) \\
T3 & $B$ & \texttt{Sigma.fst} & dependent $E\,b$ & \texttt{sigmaSection} (\texttt{noncomputable}) \\
T4 & \texttt{Bool} & \texttt{forgetSumIsLeft} & tagged \texttt{Sum} values & \texttt{sumSection} (parametric) \\
T5 & $\alpha$ & \texttt{Subtype.val} & proof data (singleton fibers) & \texttt{subtypeSection} \\
T6 & orbit quotient & \texttt{to\_orbit\_quotient} & orbit classes & \texttt{Quotient.out} (canonical) \\
T7 & $R/I$ & \texttt{forgetToIdealQuotient} & cosets mod $I$ & \texttt{Quotient.out} (canonical) \\
T8 & pushout & \texttt{forgetToPushout} & \texttt{EqvGen} classes & \texttt{Quotient.out} (canonical) \\
T9 & $\alpha$ & \texttt{Subtype.val} on \texttt{S} & membership fibers & \texttt{classifierSection} \\
T10 & localized ring & \texttt{forgetPairToLocalization} & denominator pairs & \texttt{IsLocalization.sec} \\
T11 & $A$ & \texttt{pullbackFst} & compatible $B$-legs & \texttt{pullbackSection} (\texttt{noncomputable}) \\
T12 & \texttt{Nat} & \texttt{Nat.succ} & predecessor fibers & \texttt{Nat.pred} (left inverse only) \\
\bottomrule
\end{tabular}
\caption{Side-by-side view: twelve instantiations (T12 is a MODERATE boundary case).}
\label{tab:synth}
\end{table}

\noindent
\textbf{T1} isolates \emph{collapse by an equivalence relation}; \textbf{T2} isolates
\emph{multi-route} without quotienting; \textbf{T3} is the dependent-type generalization of
T2; \textbf{T4} is the coproduct-side dual tag story; \textbf{T5} is the subtype refinement
pattern; \textbf{T6} and \textbf{T7} are quotient benchmarks (group action or ideal).
\textbf{T8}--\textbf{T11} extend the story to colimits, classifiers, localization pairs, and pullbacks.
\textbf{T12} is the explicit \textbf{MODERATE} control. For \textbf{T1}--\textbf{T11}, the
\textbf{STRONG} rubric applies where stated; \textbf{T12} documents a boundary where global
\texttt{HasRightInverse} fails.

\section{Unified validation verdict and the meaning of \textbf{STRONG}}\label{sec:verdict}

Under the project's validation rubric~\cite{EPIC015WV1}, a benchmark is rated \textbf{STRONG}
when it exports the architecture cleanly to external mathematics and yields a theorem bundle
structurally richer than the obvious baseline. Tranches \textbf{T1}--\textbf{T11} are recorded
with that ambition; \textbf{T12} is an explicit \textbf{MODERATE} / methodological control:
non-surjective collapse with nonempty fibers only on the image. \emph{Collectively}, the
tranches strengthen the claim that the pattern is not tied to a single family of examples.

\textbf{STRONG} encodes methodological transfer and corollary structure, not difficulty of
raw mathematical content. No tranche proves a new ``summit'' theorem for the \IC{} landscape;
together they support external validity of the \emph{packaging}.

\section{Place in the larger program}\label{sec:stack}

\begin{itemize}[leftmargin=2em]
  \item \textbf{\IC{} flagship (internal):} the positive-closure architecture and main
    certification / reflective splits---load-bearing inside the library.
  \item \textbf{External validation (this paper):} machine-checked demonstrations that the
    same pattern applies \emph{outside} the flagship modules, on quotients, products, dependent
    sums, coproduct discriminants, subtypes, orbit and ideal quotients, pushouts, classifiers,
    localization pairs, pullbacks, and a non-surjective control tranche.
  \item \textbf{Strata:} abstract synthesis and bridge laws~\cite{SpivackStrata}.
  \item \textbf{APS:} local-to-global backbone elsewhere in the stack. \textbf{These tranches
    do not test APS directly;} they test reusable positive-closure \emph{packaging} on
    external mathematics.
  \item \textbf{NEMS:} outer burden and interface story for the paper suite.
\end{itemize}

\section{Limitations and future work}\label{sec:limits}

Twelve finite families do not exhaust all closure or fibration patterns. Natural extensions
include: settings where sections are non-unique or only locally defined; competing enriched
layers over one collapse map; and domain-specific benchmarks outside pure type theory. Such
work would extend the \emph{evidence base} without changing the logical status of the
positive-closure architecture inside \IC{}.

\appendix

\section{Theorem indices (Lean names by module)}\label{app:index}

\subsection*{\texttt{QuotientFiberBenchmark.lean}}

{\footnotesize
\begin{longtable}{@{}>{\raggedright\arraybackslash}p{0.40\textwidth}>{\raggedright\arraybackslash}p{0.54\textwidth}@{}}
\toprule
\textbf{Lean name} & \textbf{One-line reading} \\
\midrule
\endfirsthead
\multicolumn{2}{c}{\textit{(continued)}} \\
\toprule
\textbf{Lean name} & \textbf{One-line reading} \\
\midrule
\endhead
\midrule
\multicolumn{2}{r}{\textit{continued on next page}} \\
\endfoot
\bottomrule
\endlastfoot

\texttt{BareQuotient} &
quotient type as bare certificate layer after collapse. \\

\texttt{EnrichedCarrier} &
carrier before quotienting (enriched layer). \\

\texttt{forgetToQuotient} &
projection \texttt{Quotient.mk} from carrier to quotient. \\

\texttt{QuotientFiber} &
fiber over a bare quotient point (subtype of the carrier). \\

\texttt{quotient\_fiber\_nonempty} &
each bare point has a nonempty fiber (from \texttt{exists\_rep}). \\

\texttt{forgetToQuotient\_surjective} &
\texttt{forgetToQuotient} is surjective onto the quotient. \\

\texttt{forgetToQuotient\_hasRightInverse} &
projection has a global right inverse (section exists). \\

\texttt{quotientOut\_rightInverse\_forgetToQuotient} &
\texttt{Quotient.out} is a right inverse of \texttt{forgetToQuotient}. \\

\texttt{canonicalQuotientFiberWitness} &
canonical distinguished point in each \texttt{QuotientFiber}. \\

\texttt{quotient\_fiber\_nonempty'} &
nonemptiness via \texttt{canonicalQuotientFiberWitness}. \\

\end{longtable}
}

\subsection*{\texttt{ProductProjectionFiberBenchmark.lean}}

{\footnotesize
\begin{longtable}{@{}>{\raggedright\arraybackslash}p{0.40\textwidth}>{\raggedright\arraybackslash}p{0.54\textwidth}@{}}
\toprule
\textbf{Lean name} & \textbf{One-line reading} \\
\midrule
\endfirsthead
\multicolumn{2}{c}{\textit{(continued)}} \\
\toprule
\textbf{Lean name} & \textbf{One-line reading} \\
\midrule
\endhead
\bottomrule
\endlastfoot

\texttt{forgetFirst} & first projection $X \times Y \to X$. \\
\texttt{ProductFiberAt X Y x} & fiber of \texttt{forgetFirst} over \texttt{x}. \\
\texttt{productSection} & section $x \mapsto (x, y_0)$ at fixed \texttt{y0}. \\
\texttt{product\_fiber\_nonempty} & nonempty \texttt{ProductFiberAt} given \texttt{y0}. \\
\texttt{forgetFirst\_surjective} & \texttt{forgetFirst} surjective given \texttt{y0}. \\
\texttt{forgetFirst\_hasRightInverse} & \texttt{HasRightInverse} for \texttt{forgetFirst}. \\
\texttt{productSection\_rightInverse\_forgetFirst} & \texttt{RightInverse} law. \\
\texttt{canonicalProductFiberWitness} & canonical witness in each fiber. \\
\texttt{product\_fiber\_nonempty'} & nonemptiness via canonical witness. \\

\end{longtable}
}

\subsection*{\texttt{SigmaFiberBenchmark.lean}}

{\footnotesize
\begin{longtable}{@{}>{\raggedright\arraybackslash}p{0.40\textwidth}>{\raggedright\arraybackslash}p{0.54\textwidth}@{}}
\toprule
\textbf{Lean name} & \textbf{One-line reading} \\
\midrule
\endfirsthead
\multicolumn{2}{c}{\textit{(continued)}} \\
\toprule
\textbf{Lean name} & \textbf{One-line reading} \\
\midrule
\endhead
\bottomrule
\endlastfoot

\texttt{SigmaFiber} & fiber of \texttt{Sigma.fst} over \texttt{b}. \\
\texttt{sigma\_fiber\_nonempty} & nonempty fibers from global nonemptiness. \\
\texttt{sigma\_fst\_surjective} & \texttt{Sigma.fst} surjective. \\
\texttt{sigma\_fst\_hasRightInverse} & \texttt{HasRightInverse}. \\
\texttt{chooseFiber} & \texttt{Classical.choice} in each fiber. \\
\texttt{sigmaSection} & global section of the bundle. \\
\texttt{sigmaSection\_rightInverse\_sigma\_fst} & \texttt{RightInverse} law. \\
\texttt{canonicalSigmaFiberWitness} & canonical witness in each \texttt{SigmaFiber}. \\

\end{longtable}
}

\subsection*{\texttt{SumCoproductFiberBenchmark.lean}}

{\footnotesize
\begin{longtable}{@{}>{\raggedright\arraybackslash}p{0.40\textwidth}>{\raggedright\arraybackslash}p{0.54\textwidth}@{}}
\toprule
\textbf{Lean name} & \textbf{One-line reading} \\
\midrule
\endfirsthead
\multicolumn{2}{c}{\textit{(continued)}} \\
\toprule
\textbf{Lean name} & \textbf{One-line reading} \\
\midrule
\endhead
\bottomrule
\endfoot

\texttt{forgetSumIsLeft} & discriminant collapse \texttt{Sum A B → Bool}. \\
\texttt{SumFiberAt} & fiber over a side tag. \\
\texttt{sumSection} & parametric section from \texttt{Bool}. \\
\texttt{sum\_fiber\_nonempty} & nonempty fibers (given \texttt{a0}, \texttt{b0}). \\
\texttt{forgetSumIsLeft\_surjective} & surjectivity. \\
\texttt{forgetSumIsLeft\_hasRightInverse} & \texttt{HasRightInverse}. \\
\texttt{sumSection\_rightInverse\_forgetSumIsLeft} & \texttt{RightInverse} law. \\
\texttt{canonicalSumFiberWitness} & canonical witness in each fiber. \\

\end{longtable}
}

\subsection*{\texttt{SubtypeValFiberBenchmark.lean}}

{\footnotesize
\begin{longtable}{@{}>{\raggedright\arraybackslash}p{0.40\textwidth}>{\raggedright\arraybackslash}p{0.54\textwidth}@{}}
\toprule
\textbf{Lean name} & \textbf{One-line reading} \\
\midrule
\endfirsthead
\multicolumn{2}{c}{\textit{(continued)}} \\
\toprule
\textbf{Lean name} & \textbf{One-line reading} \\
\midrule
\endhead
\bottomrule
\endfoot

\texttt{SubtypeFiberAt} & fiber of \texttt{Subtype.val} over \texttt{a}. \\
\texttt{subtype\_fiber\_nonempty} & nonempty fibers from $\forall a,\ p\ a$. \\
\texttt{subtype\_val\_surjective} & surjectivity of \texttt{Subtype.val}. \\
\texttt{subtype\_val\_hasRightInverse} & \texttt{HasRightInverse}. \\
\texttt{subtypeSection} & section sending \texttt{a} to \texttt{<a, h a>}. \\
\texttt{subtypeSection\_rightInverse\_subtypeVal} & \texttt{RightInverse} law. \\
\texttt{canonicalSubtypeFiberWitness} & canonical witness in each fiber. \\

\end{longtable}
}

\subsection*{\texttt{OrbitRelationQuotientFiberBenchmark.lean}}

{\footnotesize
\begin{longtable}{@{}>{\raggedright\arraybackslash}p{0.40\textwidth}>{\raggedright\arraybackslash}p{0.54\textwidth}@{}}
\toprule
\textbf{Lean name} & \textbf{One-line reading} \\
\midrule
\endfirsthead
\multicolumn{2}{c}{\textit{(continued)}} \\
\toprule
\textbf{Lean name} & \textbf{One-line reading} \\
\midrule
\endhead
\bottomrule
\endfoot

\texttt{same\_orbit} & orbit relation from a \texttt{MulAction}. \\
\texttt{grp\_action\_orbit\_setoid} & \texttt{Setoid} for orbit equivalence. \\
\texttt{BareOrbitQuotient} & orbit-space quotient type. \\
\texttt{to\_orbit\_quotient} & collapse map \texttt{Quotient.mk}. \\
\texttt{OrbitQuotientFiber} & fiber over an orbit point. \\
\texttt{orbit\_quotient\_fiber\_nonempty} & nonempty fibers. \\
\texttt{to\_orbit\_quotient\_surjective} & surjectivity. \\
\texttt{to\_orbit\_quotient\_hasRightInverse} & \texttt{HasRightInverse}. \\
\texttt{quotient\_out\_rightInverse\_to\_orbit\_quotient} & \texttt{Quotient.out} as \texttt{RightInverse}. \\
\texttt{canonical\_orbit\_quotient\_fiber\_witness} & canonical witness in each fiber. \\

\end{longtable}
}

\subsection*{\texttt{IdealQuotientFiberBenchmark.lean}}

{\footnotesize
\begin{longtable}{@{}>{\raggedright\arraybackslash}p{0.40\textwidth}>{\raggedright\arraybackslash}p{0.54\textwidth}@{}}
\toprule
\textbf{Lean name} & \textbf{One-line reading} \\
\midrule
\endfirsthead
\multicolumn{2}{c}{\textit{(continued)}} \\
\toprule
\textbf{Lean name} & \textbf{One-line reading} \\
\midrule
\endhead
\bottomrule
\endfoot

\texttt{forgetToIdealQuotient} & \texttt{Ideal.Quotient.mk I}. \\
\texttt{IdealQuotientFiber} & fiber over a coset in $R/I$. \\
\texttt{ideal\_quotient\_fiber\_nonempty} & nonempty fibers. \\
\texttt{forgetToIdealQuotient\_surjective} & surjectivity. \\
\texttt{forgetToIdealQuotient\_hasRightInverse} & \texttt{HasRightInverse}. \\
\texttt{idealQuotientOut\_rightInverse\_forgetToIdealQuotient} & \texttt{Quotient.out} as \texttt{RightInverse}. \\
\texttt{canonicalIdealQuotientFiberWitness} & canonical witness in each fiber. \\

\end{longtable}
}

\subsection*{\texttt{PushoutCoequalizerFiberBenchmark.lean}}

{\footnotesize
\begin{longtable}{@{}>{\raggedright\arraybackslash}p{0.40\textwidth}>{\raggedright\arraybackslash}p{0.54\textwidth}@{}}
\toprule
\textbf{Lean name} & \textbf{One-line reading} \\
\midrule
\endfirsthead
\multicolumn{2}{c}{\textit{(continued)}} \\
\toprule
\textbf{Lean name} & \textbf{One-line reading} \\
\midrule
\endhead
\bottomrule
\endfoot

\texttt{pushoutRel} & generating symmetric edges for glueing. \\
\texttt{pushoutSetoid} & \texttt{EqvGen.setoid} presentation. \\
\texttt{forgetToPushout} & collapse to the pushout quotient. \\
\texttt{PushoutFiber} & fiber over a glued class. \\
\texttt{forgetToPushout\_surjective} & surjectivity. \\
\texttt{quotientOut\_rightInverse\_forgetToPushout} & \texttt{Quotient.out} as \texttt{RightInverse}. \\
\texttt{canonicalPushoutFiberWitness} & canonical witness in each fiber. \\

\end{longtable}
}

\subsection*{\texttt{SetClassifierFiberBenchmark.lean}}

{\footnotesize
\begin{longtable}{@{}>{\raggedright\arraybackslash}p{0.40\textwidth}>{\raggedright\arraybackslash}p{0.54\textwidth}@{}}
\toprule
\textbf{Lean name} & \textbf{One-line reading} \\
\midrule
\endfirsthead
\multicolumn{2}{c}{\textit{(continued)}} \\
\toprule
\textbf{Lean name} & \textbf{One-line reading} \\
\midrule
\endhead
\bottomrule
\endfoot

\texttt{membershipClassifier} & characteristic map \texttt{x in S} to \texttt{Prop}. \\
\texttt{ClassifierFiber} & fiber of \texttt{Subtype.val} over \texttt{x}. \\
\texttt{classifier\_fiber\_nonempty\_of\_mem} & nonempty when \texttt{x} lies in \texttt{S}. \\
\texttt{subtypeVal\_surjectiveOn} & surjectivity under global membership. \\
\texttt{subtypeVal\_hasRightInverse} & \texttt{HasRightInverse} under global membership. \\
\texttt{classifierSection\_rightInverse\_subtypeVal} & \texttt{RightInverse} law. \\

\end{longtable}
}

\subsection*{\texttt{LocalizationPairFiberBenchmark.lean}}

{\footnotesize
\begin{longtable}{@{}>{\raggedright\arraybackslash}p{0.40\textwidth}>{\raggedright\arraybackslash}p{0.54\textwidth}@{}}
\toprule
\textbf{Lean name} & \textbf{One-line reading} \\
\midrule
\endfirsthead
\multicolumn{2}{c}{\textit{(continued)}} \\
\toprule
\textbf{Lean name} & \textbf{One-line reading} \\
\midrule
\endhead
\bottomrule
\endfoot

\texttt{forgetPairToLocalization} & \texttt{IsLocalization.mk'} on \texttt{R × M}. \\
\texttt{forgetPairToLocalization\_surjective} & surjectivity of the pair map. \\
\texttt{forgetPairToLocalization\_hasRightInverse} & \texttt{HasRightInverse}. \\
\texttt{sec\_rightInverse\_forgetPairToLocalization} & \texttt{IsLocalization.sec} as \texttt{RightInverse}. \\
\texttt{LocalizationPairFiber} & fiber over a localized element. \\
\texttt{canonicalLocalizationPairFiberWitness} & canonical pair witness. \\

\end{longtable}
}

\subsection*{\texttt{PullbackFiberBenchmark.lean}}

{\footnotesize
\begin{longtable}{@{}>{\raggedright\arraybackslash}p{0.40\textwidth}>{\raggedright\arraybackslash}p{0.54\textwidth}@{}}
\toprule
\textbf{Lean name} & \textbf{One-line reading} \\
\midrule
\endfirsthead
\multicolumn{2}{c}{\textit{(continued)}} \\
\toprule
\textbf{Lean name} & \textbf{One-line reading} \\
\midrule
\endhead
\bottomrule
\endfoot

\texttt{PullbackType} & pullback as subtype of \texttt{A × B}. \\
\texttt{pullbackFst} & first projection (renamed; avoids clash with T2). \\
\texttt{pullbackFst\_surjective} & surjectivity under overlap hypothesis. \\
\texttt{pullbackFst\_hasRightInverse} & \texttt{HasRightInverse}. \\
\texttt{pullbackSection\_rightInverse\_pullbackFst} & \texttt{RightInverse} for chosen section. \\

\end{longtable}
}

\subsection*{\texttt{NonSurjectiveSuccFiberBenchmark.lean}}

{\footnotesize
\begin{longtable}{@{}>{\raggedright\arraybackslash}p{0.40\textwidth}>{\raggedright\arraybackslash}p{0.54\textwidth}@{}}
\toprule
\textbf{Lean name} & \textbf{One-line reading} \\
\midrule
\endfirsthead
\multicolumn{2}{c}{\textit{(continued)}} \\
\toprule
\textbf{Lean name} & \textbf{One-line reading} \\
\midrule
\endhead
\bottomrule
\endfoot

\texttt{succFiber} & fiber of \texttt{Nat.succ}. \\
\texttt{succFiber\_nonempty\_iff} & nonempty iff not zero. \\
\texttt{nat\_succ\_not\_surjective} & no global section onto \texttt{Nat}. \\
\texttt{nat\_succ\_not\_hasRightInverse} & obstruction to \texttt{HasRightInverse}. \\
\texttt{pred\_leftInverse\_succ} & \texttt{Nat.pred} as left inverse of \texttt{succ}. \\

\end{longtable}
}

\section{A workflow for reusing this pattern}\label{app:recipe}

The benchmarks above do \emph{not} constitute a single new lemma or tactic that resolves
goals that were previously out of reach. They exhibit a \textbf{reusable organization}:
choose a surjection or collapse map $\pi : E \to B$, describe the fibers, state a minimal
existence theorem as baseline, then---when it fits---strengthen the interface with surjectivity
of $\pi$, existence of a (possibly classical) section, and canonical choices in fibers
(Section~\ref{sec:method}). This appendix gives a step-by-step way to apply that organization
to another problem, including problems that remain partly or fully open.

\paragraph{Scope.}
The workflow clarifies \emph{where} a forgetful map and fiber structure appear, and whether a
global right inverse is even available. It does not replace domain mathematics: restating a
conjecture in fiber language rarely proves it. New ideas still come from the field; this
pattern helps one see which lemmas would sit at the boundary between a coarse certificate and
the data that varies above it.

\paragraph{Step 1: isolate the baseline theorem.}
Write the statement a working mathematician would accept as the minimal ``obvious'' layer:
existence of representatives, existence of preimages under a surjection, or other standard
representability. That statement is the floor; anything layered on top should imply it, not
replace it.

\paragraph{Step 2: identify $B$, $E$, and $\pi$.}
Fix the set or type $B$ of coarse certificates, a type $E$ of richer structure, and a map
$\pi : E \to B$ whose fibers describe what is not determined by the certificate alone. Use
Table~\ref{tab:synth} and Section~\ref{sec:synth} as a menu: quotients and glueing, products
and dependent sums, tags and subtypes, group or ideal quotients, localization via pairs,
pullbacks, or---when $\pi$ is not onto---the phenomenon illustrated by successor on the natural
numbers (nonempty fibers only on the image, partial inverses instead of a global section).

\paragraph{Step 3: formalize the baseline in Lean.}
Prove or import the minimal facts in Mathlib style. If the aim is to check portability of the
pattern \emph{without} leaning on a large custom library, keep the formalization in a small
standalone module that depends only on Mathlib (as in the benchmark files accompanying this
paper).

\paragraph{Step 4: strengthen when possible.}
When $\pi$ is surjective, it is often natural to record: nonempty fibers, surjectivity of
$\pi$, existence of a right inverse in the sense of \texttt{HasRightInverse}, a concrete
\texttt{RightInverse}, and a canonical element of each fiber (Section~\ref{sec:method}). When
some of this fails, the failure is itself information: for example, there may be no global
section, only restriction to a subdomain or a one-sided inverse (as with predecessor and
successor). State that boundary case explicitly rather than forcing a full section bundle.

\paragraph{Step 5: open problems and partial progress.}
If the main theorem is still open, use the same decomposition to list separate conjectures:
what belongs to the certificate layer, what belongs to fiber data, what would be a section or
choice principle. Partial results can attach cleanly to one of these layers. Progress on a
hard problem may then consist of proving a lemma at that interface---even when the headline
claim remains open---rather than of relabeling the conjecture.

\begin{thebibliography}{99}

\bibitem{SpivackIC}
N.~Spivack.
\newblock \emph{Canonical Certification Does Not Exhaust Reflective Structure: Reflective Split, Strict Refinement, and Fiber Structure in Infinity Compression}.
\newblock Manuscript, 2026.
\newblock Companion \LaTeX{} source: \texttt{papers/Canonical\_Certification\_and\_Enriched\_Reflective\_Realization/}.

\bibitem{SpivackStrata}
N.~Spivack.
\newblock \emph{Strata: Abstract Synthesis and Bridge Laws (Infinity Compression)}.
\newblock Companion manuscript in preparation, 2026.

\bibitem{EPIC015WV1}
Infinity Compression contributors.
\newblock \emph{EPIC\_015\_WV1 --- External Proof Architecture Validation} (specification and validation verdict notes).
\newblock \texttt{infinity-compression} repository, \texttt{specs/} tree, 2026.

\bibitem{Lean4}
L.~de Moura and S.~Ullrich.
\newblock The Lean 4 theorem prover and programming language.
\newblock In \emph{Automated Deduction --- CADE 28}, volume 12699 of \emph{Lecture Notes in Computer Science}, pages 625--635. Springer, 2021.

\bibitem{Mathlib}
The Mathlib community.
\newblock The Lean mathematical library.
\newblock In \emph{Proceedings of the 9th ACM SIGPLAN International Conference on Certified Programs and Proofs (CPP)}, pages 367--381, 2020.
\newblock \textsc{doi}: \href{https://doi.org/10.1145/3372885.3373824}{10.1145/3372885.3373824}.
\newblock Mathlib~4 repository: \url{https://github.com/leanprover-community/mathlib4}.

\bibitem{MacLane}
S.~Mac Lane.
\newblock \emph{Categories for the Working Mathematician}, chapter I (categories of algebraic structures; universal constructions motivating quotients).
\newblock Springer, 2nd edition, 1998.

\end{thebibliography}

\end{document}
